\documentclass[aspectratio=169,11pt]{beamer}

\usepackage[utf8]{inputenc}
\usepackage[T1]{fontenc}
\usepackage[spanish]{babel}
\usepackage{graphicx}
\usepackage{booktabs}
\usepackage{amsmath}
\usepackage{hyperref}

\usetheme{Madrid}
\usecolortheme{default}
\setbeamertemplate{navigation symbols}{}
\setbeamertemplate{caption}[numbered]

\title[ACP calidad del agua]{Análisis de Componentes Principales aplicado a calidad del agua del Río Cauca}
\subtitle{Taller 1. Análisis Multivariado}
\author[J. Tapia, J. Jaramillo, L. Meza]{
JHONATAN ANDRES TAPIA CORDOBA\\
JORGE ANDRES JARAMILLO NEME\\
LUIS FERNANDO MEZA RAMREZ
}
\institute{Maestría en Inteligencia Artificial y Ciencia de Datos}
\date{Mayo 2 de 2026}

\newcommand{\img}[2]{
  \begin{center}
    \includegraphics[width=#2\textwidth,height=0.58\textheight,keepaspectratio]{\detokenize{#1}}
  \end{center}
}

\begin{document}

\begin{frame}
  \titlepage
\end{frame}

\begin{frame}{Objetivo y alcance}
  \begin{itemize}
    \item Aplicar ACP sobre variables físico-químicas de calidad del agua del Río Cauca.
    \item Resumir la información multivariada en dimensiones interpretables.
    \item Identificar patrones, observaciones atípicas y posibilidad de construir un índice.
  \end{itemize}

  \vspace{0.25cm}
  \textbf{Idea central:} la calidad del agua no se resume completamente en una sola variable ni en un solo eje factorial.
\end{frame}

\begin{frame}{Datos y preparación}
  \begin{columns}[T]
    \begin{column}{0.48\textwidth}
      \begin{itemize}
        \item Base original: \textbf{2368 registros y 56 variables}.
        \item Limpieza de nombres y formatos numéricos.
        \item Variables sin información útil fueron descartadas.
        \item Individuo: combinación \texttt{anio\_estacion}.
      \end{itemize}
    \end{column}
    \begin{column}{0.48\textwidth}
      \begin{itemize}
        \item Variables activas: turbiedad, SST, DBO, DQO, OD, conductividad y fósforo.
        \item Transformación \texttt{log1p} por asimetría positiva severa.
        \item Estandarización por unidades diferentes.
        \item Matriz final: \textbf{1475 observaciones} (\textbf{62,29 \%}).
      \end{itemize}
    \end{column}
  \end{columns}
\end{frame}

\begin{frame}{Exploración y correlaciones}
  \begin{itemize}
    \item Alta heterogeneidad, especialmente en DBO, fósforo total y oxígeno disuelto.
    \item Turbiedad y sólidos suspendidos totales: \textbf{r = 0,80}.
    \item DBO y fósforo total: \textbf{r = 0,60}.
    \item Oxígeno disuelto y conductividad eléctrica: \textbf{r = -0,54}.
  \end{itemize}

  \img{Graficos de resultados/Matriz de correlacion: Espacion de variable.png}{0.88}
\end{frame}

\begin{frame}{Varianza explicada}
  \begin{columns}[T]
    \begin{column}{0.42\textwidth}
      \begin{itemize}
        \item Criterio de Kaiser: \textbf{3 componentes}.
        \item Dim 1: \textbf{31,57 \%}.
        \item Dim 2: \textbf{25,43 \%}.
        \item Dim 3: \textbf{20,15 \%}.
        \item Acumulado: \textbf{77,15 \%}.
        \item Plano 1-2: \textbf{57,00 \%}.
      \end{itemize}
    \end{column}
    \begin{column}{0.58\textwidth}
      \img{Graficos de resultados/Sedimentacion de varianza.png}{1.0}
    \end{column}
  \end{columns}
\end{frame}

\begin{frame}{Interpretación de los ejes}
  \begin{itemize}
    \item \textbf{Dim 1:} dominada por sólidos suspendidos, turbiedad, DBO, DQO y fósforo total.
    \item Interpretación: \textbf{gradiente de carga contaminante, orgánica y particulada}.
    \item \textbf{Dim 2:} dominada por oxígeno disuelto y conductividad eléctrica en sentidos opuestos.
    \item Interpretación: \textbf{contraste entre oxigenación e influencia iónica/conductiva}.
  \end{itemize}

  \img{Graficos de resultados/Circulos de correlaciones Plano 1-2.png}{0.86}
\end{frame}

\begin{frame}{Individuos, biplot y atípicos}
  \begin{columns}[T]
    \begin{column}{0.52\textwidth}
      \img{Graficos de resultados/Relacion Dual individuos - variables.png}{1.0}
    \end{column}
    \begin{column}{0.48\textwidth}
      \small
      \begin{itemize}
        \item El biplot integra individuos y variables.
        \item Puntos alejados del origen tienen perfiles diferenciados.
        \item Posibles atípicos factoriales:
          \begin{itemize}
            \item 1993\_JUANCHITO.
            \item 2010\_MEDIACANOA.
            \item 2010\_PASO DEL COMERCIO.
            \item 2010\_PUERTO ISAACS.
            \item 2010\_JUANCHITO.
          \end{itemize}
      \end{itemize}
      No se eliminan automáticamente: pueden ser eventos ambientales reales.
    \end{column}
  \end{columns}
\end{frame}

\begin{frame}{Contribuciones y cosenos}
  \begin{columns}[T]
    \begin{column}{0.5\textwidth}
      \textbf{Dim 1: carga contaminante}
      \begin{itemize}
        \item Sólidos suspendidos: \textbf{32,36 \%}.
        \item Turbiedad: \textbf{23,45 \%}.
        \item DBO: \textbf{15,71 \%}.
      \end{itemize}
    \end{column}
    \begin{column}{0.5\textwidth}
      \textbf{Dim 2: oxigenación/conductividad}
      \begin{itemize}
        \item Oxígeno disuelto: \textbf{41,02 \%}.
        \item Conductividad eléctrica: \textbf{32,98 \%}.
      \end{itemize}
    \end{column}
  \end{columns}

  \vspace{0.3cm}
  \textbf{Lectura:} los cosenos cuadrados confirman que Dim 1 y Dim 2 representan fenómenos ambientales distintos y complementarios.
\end{frame}

\begin{frame}{Índice parcial}
  \begin{itemize}
    \item Sí es posible construir un índice, pero como \textbf{índice parcial de carga contaminante/particulada}.
    \item No debe interpretarse como índice global de calidad del agua.
    \item PC1 explica \textbf{31,57 \%}; PC2 conserva información relevante sobre oxígeno disuelto y conductividad.
  \end{itemize}

  \[
    \text{Indice\_PC1\_0\_100} =
    100 \times \frac{PC1 - \min(PC1)}{\max(PC1) - \min(PC1)}
  \]

  \small
  Mayores valores: 2010\_LA VICTORIA, 1994\_PASO DEL COMERCIO, 2010\_ANACARO, 2010\_MEDIACANOA y 2010\_PASO DEL COMERCIO.\\
  Mayores promedios por estación: ANACARO, PUENTE LA VIRGINIA, LA VICTORIA, VIJES y PUENTE GUAYABAL.
\end{frame}

\begin{frame}{Conclusiones}
  \begin{enumerate}
    \item El ACP resumió siete variables físico-químicas en tres componentes que explican \textbf{77,15 \%}.
    \item La Dim 1 representa carga contaminante y particulada.
    \item La Dim 2 complementa la lectura mediante oxígeno disuelto y conductividad eléctrica.
    \item Se identificaron observaciones extremas potencialmente relevantes.
    \item El índice construido es útil para ordenar observaciones y estaciones, pero es \textbf{parcial}.
  \end{enumerate}

  \vspace{0.25cm}
  \textbf{Mensaje final:} la calidad del agua del Río Cauca debe leerse conjuntamente desde carga contaminante y condiciones de oxigenación/conductividad.
\end{frame}

\begin{frame}{Preguntas probables}
  \small
  \textbf{¿Por qué log1p?} Por asimetría positiva severa y valores extremos.\\[0.22cm]
  \textbf{¿Por qué estandarizar?} Porque las variables estaban en unidades distintas.\\[0.22cm]
  \textbf{¿Por qué tres componentes?} Porque tres valores propios fueron mayores que 1 y explicaron 77,15 \%.\\[0.22cm]
  \textbf{¿Por qué el índice no es global?} Porque PC2 contiene información ambiental relevante no resumida por PC1.\\[0.22cm]
  \textbf{¿Qué significa un atípico en ACP?} Una observación alejada del centro factorial; no necesariamente un error.
\end{frame}

\end{document}
